\documentclass[11pt,a4paper]{article}
\usepackage[left=2.2cm,right=2.2cm,top=1.8cm,bottom=1.8cm]{geometry}
\usepackage[T1]{fontenc}
\usepackage{lmodern}
\usepackage[ngerman]{babel}
\usepackage[hidelinks]{hyperref}
\usepackage{parskip}
\pagestyle{empty}

\begin{document}

\begin{flushright}
Kostiantyn Pysanyi\\
Leipzig, Germany\\
\href{mailto:kostiantyn.pysanyi@gmail.com}{kostiantyn.pysanyi@gmail.com}\\
+49 1525 8447880
\end{flushright}

WetterOnline Meteorologische Dienstleistungen GmbH\\
Recruiting-Team\\
Karl-Legien-Straße 194a\\
53117 Bonn

\begin{flushright}
30. August 2026
\end{flushright}

\textbf{Bewerbung als Machine Learning Engineer (w/m/d)}

Sehr geehrtes Recruiting-Team,

ich bewerbe mich als Machine Learning Engineer, weil mich die Verbindung aus Computer Vision, großen raumzeitlichen Bilddaten und physikalisch geprägten Fragestellungen besonders anspricht. Mich interessiert besonders, wie neuronale Netze robuste Strukturen in komplexen Radar- und Satellitendaten erkennen und verständlich darstellen können. Für diese Arbeit bringe ich Erfahrung in der Entwicklung, Evaluation und produktiven Umsetzung von Bildanalysemodellen sowie einen fachlichen Hintergrund in Data Science und Ingenieurwissenschaften mit.

Bei RAYLYTIC entwickelte und trainierte ich produktionsreife Modelle für Segmentierung, Detektion und Klassifikation medizinischer Bilddaten. Je nach Aufgabe und Metrik verbesserte ich die Modellgüte um 20--50\% gegenüber den Baselines und verkürzte die Trainingszeit durch optimierte Datenverarbeitung um den Faktor acht. Außerdem baute ich eine 2D--3D-Registrierungspipeline für die Implantatlagenschätzung von der Bild- und Mesh-Verarbeitung bis zum GPU-beschleunigten Deployment auf. In einem weiteren Forschungsprojekt entwickelte ich 3D-Instanzsegmentierung und Registrierung für volumetrische Lichtblattmikroskopie. Während meiner Tätigkeit als wissenschaftlicher Mitarbeiter sammelte ich zudem praktische Erfahrung mit verteilten Trainingsverfahren auf dem HPC-Cluster der Hochschule und brachte moderne Trainingsframeworks wie PyTorch Lightning in die Arbeitsgruppe ein. Dadurch bin ich es gewohnt, komplexe Bilddaten systematisch aufzubereiten, Modelle reproduzierbar zu trainieren und ihre Ergebnisse anhand klarer Metriken zu verbessern.

An WetterOnline reizt mich besonders, an wissenschaftlich anspruchsvollen Problemen zu arbeiten, die zugleich gesellschaftliche Bedeutung haben. Die Analyse von Wetterdaten hilft nicht nur bei präziseren Vorhersagen, sondern auch dabei, die Prozesse besser zu verstehen, die Wetter und Klima prägen; gerade angesichts des Klimawandels werden verlässliche Wetterinformationen künftig noch wichtiger. Ich möchte meine Erfahrung in Computer Vision und robusten ML-Systemen einbringen, um aus komplexen Radar- und Satellitendaten belastbare und verständliche Ergebnisse zu gewinnen.

Ich freue mich darauf, meine Erfahrungen und meine Motivation in einem persönlichen Gespräch näher zu erläutern.

Mit freundlichen Grüßen\\[0.8em]
Kostiantyn Pysanyi

\end{document}
